% Round 16. Explicit actual activation of historical/published positivity.
% Mathematical source independently PASS by higher_trace at SHA
% 203455da36d1200359656dd98e77fdd565fab9e982b000e41096859568ef5ae5.
% This activation changes only provenance comments. No new positivity mechanism.
\section{Positive fourfold components: an explicit actual activation}
\label{r16pos:section}

\paragraph{Provenance.}
The Ancona application and the projective-bundle signature calculations
below already appear in the historical Round 7 source
\path{history_01_07/bb_round7/proofs/threefold_fourfold_examples.tex}.
Its supplied-arithmetic-cover consequence retained a human-review flag.
Round 15 subsequently proved the multiplicity-weighted isotropy and
its exact-support actual bridge with the revised pure-operational
interfaces. The present section records the precise local consequence
of those accepted inputs, adds the full fourfold conclusion using
the new one-cycle theorem, and states its original-field dominant
witness scope. It does not claim discovery of the positivity
argument, a new theorem of Hodge type, or existence of such a model
for an arbitrary fourfold.

\begin{definition}[Primitive-positive numerical fourfold]
\label{r16pos:definition}
For a smooth projective geometrically integral fourfold $D/k$ over
a finite field and an ample rational divisor class $H$, put
\[
 P_H^2(D)=\ker\bigl(H:N^2(D)_{\mathbf Q}
                              \longrightarrow N^3(D)_{\mathbf Q}\bigr).
 \tag{PP16.1}\label{r16pos:primitive}
\]
The pair $(D,H)$ is primitive-positive if
$\deg_k(xy)$ is positive definite on $P_H^2(D)_{\mathbf R}$.
For a smooth integral $D$ with finite constant field
$k_D=H^0(D,\mathcal O_D)$, the same definition uses native
$k$-degrees, equivalently $k_D$-degrees multiplied by $[k_D:k]$.
\end{definition}

\begin{lemma}[Published abelian positivity and elementary component types]
\label{r16pos:types}
Every ample class gives a primitive-positive pair for any of the
following smooth projective geometrically integral fourfolds over
a finite field:
\begin{enumerate}
\item an abelian fourfold;
\item $\mathbf P_B(E)$ with $\dim B=3$ and $\operatorname{rank}E=2$;
\item $\mathbf P_S(E)$ with $\dim S=2$ and $\operatorname{rank}E=3$;
\item a projective $3$-space bundle over a smooth projective curve,
      or $\mathbf P^4$.
\end{enumerate}
Furthermore, if $k'/k$ is finite and
$f:D'\to D_{k'}$ is finite surjective with $D'$ of one of these
types, then $(D,H)$ is primitive-positive for every ample $H$.
More generally the last assertion holds whenever
$(D',f^*H_{k'})$ is primitive-positive.
\end{lemma}

\begin{proof}
The first assertion is exactly Ancona,
\emph{Standard conjectures for abelian fourfolds},
Invent. Math. \textbf{223} (2021), 149--212,
Theorems 1.3--1.4, with the numerical primitive definition in
Conjecture 1.2. The primary author version is
\url{https://arxiv.org/pdf/1806.03216}, version 3;
the DOI is \texttt{10.1007/s00222-020-00990-7}.
In dimension four and codimension two the Hodge sign is $+1$
and there is no remaining Lefschetz factor. Replacing an ample
line bundle by a very ample power leaves its primitive kernel
unchanged. No assertion that numerical and homological equivalence
coincide for a chosen $\ell$ is used.

Here are the elementary signature checks for the other entries,
also retaining their historical provenance. The numerical middle
form is nondegenerate, and on $H N^1(D)$ has signature
$(1,\rho_1(D)-1)$ by divisor Hodge index. Its orthogonal complement
is $P_H^2(D)$ by perfect numerical pairing with $N^3(D)$.
It is therefore enough to have middle signature
$(\rho_2(D)-\rho_1(D)+1,\rho_1(D)-1)$.

For a rank-two bundle on $B$, numerical projective decomposition
writes a middle class as $(\alpha,\beta)\in N^2(B)\oplus N^1(B)$.
Its pairing is
\[
 \deg_B\bigl(\alpha\beta'+\alpha'\beta+c_1(E)\beta\beta'\bigr).
\]
The change $\alpha\mapsto\alpha+\frac12c_1(E)\beta$ makes it
the hyperbolic form on $N^1(B)$ and its numerical dual.
If $\rho=\dim N^1(B)$, then
$\rho_1(D)=\rho+1$, $\rho_2(D)=2\rho$, and the signature is
$(\rho,\rho)$, as required. Numerical decomposition is justified
by the Chow projective formula: its coefficient extraction maps
are combinations of push, pull, and intersection, all of which
preserve numerical equivalence by the degree projection formula.

For a rank-three bundle on $S$, choose a degree-one rational
zero-cycle class $u$ and put $P=\pi^*u$, $\xi=c_1(\mathcal O(1))$.
The block $\xi\pi^*N^1(S)$ has the surface divisor pairing.
Replacing $\xi\pi^*D$ by
$\xi\pi^*D-\deg_S(c_1(E)D)P$ makes it orthogonal to
the plane spanned by $P,\xi^2$. That plane has matrix
\[
 \begin{pmatrix}0&1\\1&\deg_S(c_1(E)^2-c_2(E))\end{pmatrix}
\]
and signature $(1,1)$. The total signature is $(2,\rho_1(S))$,
while $\rho_1(D)=\rho_1(S)+1$ and $\rho_2(D)=\rho_1(S)+2$.
Finally the projective-space bundle over a curve has middle
hyperbolic plane and $\rho_1=\rho_2=2$; projective four-space
has $\rho_1=\rho_2=1$ and middle form positive. Their primitive
spaces are zero. These computations hold for every ample $H$
because the middle form's signature is independent of $H$.

For the finite-cover assertion, pullback on numerical groups
is injective: proper push composed with finite flat pullback
is multiplication by the full degree. A finite map between the
smooth equidimensional varieties here is flat; its generic
degree includes inseparability if present. Finite pullback
preserves ampleness and sends a primitive class to a primitive
class. Constant extension is also injective by its degree
projection formula. For a nonzero $x\in P_H^2(D)_{\mathbf R}$,
\[
  \deg_{k'}\bigl((f^*x_{k'})^2\bigr)
       =\deg(f)\deg_k(x^2)>0.
 \tag{PP16.2}\label{r16pos:finite-degree}
\]
Here degree over $k'$ of a base-changed zero-cycle equals its
original degree over $k$; viewing that same cycle over $k$
would additionally multiply it by $[k':k]$. Both conventions
give the same positive sign. This proves the asserted
positive definiteness without changing its numerical category.
\end{proof}

\begin{theorem}[Full actual G from positive fourfold components]
\label{r16pos:local}
Let $K/\mathbf Q_p$ be finite and let $X/K$ be a smooth projective
geometrically integral fourfold. Suppose that there are a finite
extension $F/K$, a smooth projective geometrically integral
fourfold $V/F$, and a projective dominant morphism
$f:V\to X_F$. Suppose $V$ has an integral regular flat projective
model $\mathcal V/F^\circ$ with
\[
             [\mathcal V_{k_F}]=\sum_i m_iD_i,\qquad m_i>0,
 \tag{PP16.3}\label{r16pos:model}
\]
whose reduced support is strict normal crossings and all closed
strata are smooth over $k_F$. Let $H$ be one relatively ample
line bundle on this model, and assume every $(D_i,H|_{D_i})$
is primitive-positive with native residue degrees.

Then every geometrically homologically trivial rational class
on $X$, in every positive codimension, has zero actual obstruction.
In particular, for every specified such codimension-$r$ cycle $Z$
there is an ordinary actual form on its exact complement,
\[
 \gamma_Z\in B^{r-1,r-1}(X^{\rm an}\setminus|Z|^{\rm an}),
 \qquad
 \int\gamma_Z\wedge D\eta=-\int_Z\eta
 \quad(\eta\in B_c^{4-r,4-r}(X^{\rm an})).
 \tag{PP16.4}\label{r16pos:G}
\]
The components may, in particular, be any mixture of the types
and finite covers specified in Lemma~\ref{r16pos:types}.
There is no bound on their number or on their primitive ranks.
The conclusion does not require the generic fourfold itself
to be abelian or to have an abelian Chow motive.
\end{theorem}

\begin{proof}
For $z\in\operatorname{CH}^2(X)_{\mathbf Q,\mathrm{geom\,hom}}$
the class $f^*z_F$ is geometrically homologically trivial on $V$.
Apply the frozen canonical residual construction over $F$,
\path{round15/actual_local/canonical_primitive_numerical_residual.tex},
Theorem \texttt{r15res:canonical-residual}. Its residual lies in
$\bigoplus_iP_H^2(D_i)$. The weighted isotropy theorem
\texttt{r15iso:total-isotropy} in
\path{round15/actual_local/weighted_primitive_isotropy.tex} gives
\[
              \sum_i m_i\deg_{k_F}(R_i^2)=0.
 \tag{PP16.5}\label{r16pos:weighted-zero}
\]
Each summand is nonnegative and each $m_i$ is strictly positive.
Positive definiteness forces every $R_i=0$. For a component
with constant field $k_i$, its native degree is
$[k_i:k_F]\deg_{k_i}$, another positive factor; it is not
silently replaced by degree one. Thus the frozen residual
theorem supplies a rational model Chow lift with all component
numerical restrictions zero. Its pure operational $K_0$
representative, from frozen R12 \texttt{r12n:K0}, has
degree-two Chern degree zero on every vertical two-dimensional
test cycle, and generic degree-two character exactly $f^*z_F$.

Now apply Corollary~\ref{r16G:numerical-witness} to the composite
$V\to X_F\to X$. This uses that model witness in the original
$K$-theory. Viewed over $K^\circ$, the special fiber has the
same reduced components, its multiplicities acquire the
ramification factor, and vertical degree acquires the residue
factor $[k_F:k_K]$. Zero numerical degrees stay zero. The
actual model-Chern detector therefore gives the upstairs
obstruction zero in the original theory; the independently
proved dominant Bott--Chern injection gives
$\mathfrak o_{X,K}^2(z)=0$. No alteration is presumed finite
off $Z$, and no Green form is pushed along $f$.

Codimension three is the new all-dimensional one-cycle theorem
\ref{r16G:one-cycles}. Codimensions one and four are the
frozen divisor and geometric degree-zero endpoint results.
These exhaust the positive codimensions on a fourfold.
For each specified final cycle, take its own frozen curved
Green pair and subtract the global actual primitive supplied
by obstruction vanishing. This gives \eqref{r16pos:G} on
the exact original support complement and on all original
actual tests, as in Proposition~\ref{r16G:dominant-witness}.
\end{proof}

\begin{corollary}[The corresponding full complementary fourfold height]
\label{r16pos:height}
Let $X/F$ be a smooth projective geometrically integral fourfold
over a number field. Assume that at every finite place not
covered by the frozen good-model theorem the hypotheses of
Theorem~\ref{r16pos:local} are supplied over $F_v$.
Then the independently reviewed height theorem gives the
canonical normalized real pairing on the full groups
\[
 \operatorname{CH}^2(X)_{\mathbf Q,\mathrm{geom\,hom}}
       \times
 \operatorname{CH}^3(X)_{\mathbf Q,\mathrm{geom\,hom}}.
\]
The assertion includes all classes in those groups and requires
neither a global regular model nor a global model for the
local auxiliary extensions. It does not assert that all
fourfolds meet these local component conditions.
\end{corollary}

\begin{proof}
The preceding theorem and the stated good-model inputs put
both full groups in the local obstruction-zero domain at
every finite place. Apply Theorem~\ref{r16h:global-height}.
Its native local normalization, finite-place finiteness and
principal transformation law have been separately proved.
\end{proof}

\paragraph{Remaining distinction.}
The positivity hypothesis is a theorem for the explicitly
listed component classes; it is open for a general smooth
projective fourfold over a finite field. An arbitrary unpolarized
Chow-motive factorization through abelian varieties is not used
to transfer a positive primitive form. Such a transfer would
need compatibility with the primitive pairing, not merely an
isomorphism of underlying motive objects. Nor does a nonzero
residual on some other model disprove strong G.
